% iscafig.tex — paste once into the IEEEtran preamble, before \begin{document}.
% Every figure in this directory uses these styles. Do not restyle a single figure.
\ProvidesFile{iscafig.tex}[2026/10/19 uPyCore ISCA figure style]
\RequirePackage{textcomp}
\RequirePackage{tikz}
\usetikzlibrary{arrows.meta,positioning,fit,calc,backgrounds,decorations.pathreplacing}

\newcommand{\upycore}{\textmu PyCore}

\definecolor{iscaBlue}{RGB}{27,74,122}
\definecolor{iscaBlueFill}{RGB}{226,236,246}
\definecolor{iscaGreen}{RGB}{24,92,48}
\definecolor{iscaGreenFill}{RGB}{226,240,226}
\definecolor{iscaAmber}{RGB}{140,84,0}
\definecolor{iscaAmberFill}{RGB}{255,236,208}
\definecolor{iscaInk}{RGB}{35,40,46}
\definecolor{iscaSlateFill}{RGB}{241,243,245}
\definecolor{iscaRose}{RGB}{140,42,42}
\definecolor{iscaRoseFill}{RGB}{250,228,226}
\definecolor{iscaFree}{RGB}{120,126,132}

\tikzset{
  iscafig/.style={
    font=\footnotesize,
    line cap=round,
    line join=round,
    >={Stealth[length=1.7mm,width=1.25mm]},
    every node/.style={inner sep=2pt}
  },
  unit/.style={
    draw=iscaInk,
    line width=0.55pt,
    align=center,
    inner sep=3.5pt,
    font=\footnotesize,
    fill=white
  },
  core/.style={unit, fill=iscaBlueFill, draw=iscaBlue},
  mem/.style={unit, fill=iscaGreenFill, draw=iscaGreen},
  accel/.style={unit, fill=iscaAmberFill, draw=iscaAmber},
  imm/.style={unit, fill=iscaSlateFill, draw=iscaInk},
  hot/.style={unit, fill=iscaRoseFill, draw=iscaRose},
  fld/.style={
    draw=iscaInk,
    line width=0.45pt,
    align=center,
    font=\footnotesize,
    inner sep=1.5pt,
    outer sep=0pt,
    minimum height=4.7mm
  },
  fldcore/.style={fld, fill=iscaBlueFill, draw=iscaBlue},
  fldmem/.style={fld, fill=iscaGreenFill, draw=iscaGreen},
  fldacc/.style={fld, fill=iscaAmberFill, draw=iscaAmber},
  fldimm/.style={fld, fill=iscaSlateFill, draw=iscaInk},
  fldhot/.style={fld, fill=iscaRoseFill, draw=iscaRose},
  arr/.style={-{Stealth[length=1.7mm]}, draw=iscaInk, line width=0.55pt},
  thinarr/.style={-{Stealth[length=1.4mm]}, draw=iscaInk, line width=0.4pt},
  off/.style={font=\scriptsize\ttfamily, inner sep=1pt, text=iscaInk},
  lbl/.style={font=\scriptsize, align=center, text=iscaInk},
  slot/.style={
    draw=iscaInk,
    line width=0.4pt,
    minimum width=9.2mm,
    minimum height=6.4mm,
    inner sep=0pt,
    outer sep=0pt,
    font=\footnotesize,
    align=center
  },
  occ/.style={slot, fill=iscaBlueFill, draw=iscaBlue},
  lnk/.style={slot, fill=iscaRoseFill, draw=iscaRose},
  fre/.style={slot, fill=white, draw=iscaFree, dashed},
  idx/.style={font=\scriptsize\ttfamily, inner sep=0pt, text=black!70},
  grp/.style={draw=iscaInk, line width=0.6pt, inner sep=5pt},
  port/.style={
    draw=iscaInk,
    line width=0.7pt,
    fill=white,
    font=\scriptsize,
    align=center,
    inner sep=2pt,
    minimum height=4.2mm
  }
}

% \rfslot{prefix}{physical index}{style}{text}
% Eight columns. Row 0 is physical 0--7, row 1 is physical 8--15.
\newcommand{\rfslot}[4]{%
  \pgfmathtruncatemacro{\rfcol}{mod(#2,8)}%
  \pgfmathtruncatemacro{\rfrow}{floor(#2/8)}%
  \node[slot, #3] (#1#2) at ({\rfcol*1.02}, {-\rfrow*1.22}) {#4};%
}
\newcommand{\rfnums}{%
  \foreach \rfi in {0,...,15} {%
    \pgfmathtruncatemacro{\rfcol}{mod(\rfi,8)}%
    \pgfmathtruncatemacro{\rfrow}{floor(\rfi/8)}%
    \node[idx, anchor=south] at ({\rfcol*1.02}, {-\rfrow*1.22+0.36}) {\rfi};%
  }%
}

% \rec{y}{offset}{style}{width}{body}
\newcommand{\rec}[5]{%
  \node[off, anchor=east] at (-0.08,#1) {#2};%
  \node[fld, #3, anchor=west, minimum width=#4] at (0,#1) {#5};%
}
